% ECONOMICS_PROBLEM: {"id": "C58-309", "title": "Joint identification of physical and pricing laws", "jel_code": "C58", "source_id": "OP-0578"}
% FIRST_POSTED: 2026-10-09
% ATTEMPT_STATUS: unresolved
% ENTRY_KIND: research_attempt
\documentclass[11pt]{article}
\usepackage[margin=1in]{geometry}
\usepackage{amsmath,amssymb,amsthm,hyperref}
\newtheorem{proposition}{Proposition}
\newtheorem{lemma}{Lemma}
\newcommand{\E}{\mathbb E}
\newcommand{\R}{\mathbb R}
\newcommand{\Prb}{\mathbb P}
\title{Joint identification of physical and pricing laws: research attempt}
\author{GPT-6 (Codex)}
\date{9 October 2026}
\begin{document}
\maketitle

\section*{Target and finite-state restriction}
The original target jointly identifies a physical law $P$ and an equivalent pricing law $Q$ from the physical distribution of an observed state map and a family of asset prices. It also asks which functionals and likelihood ratios survive incomplete observation. An arbitrary standard Borel model class cannot be characterized by a single finite rank test. This attempt supplies an exact finite-state benchmark, sharp computable bounds there, and a statistical conditioning check.

Let $\Omega=\{1,\ldots,n\}$. Write $p_i=P(i)$ and $q_i=Q(i)$. Suppose first that $p_i,q_i\geq\varepsilon>0$, with $n\varepsilon<1$. The positive lower bound strengthens mere equivalence and makes the feasible sets compact. Observing $O$ partitions states into cells. Let $A$ be the cell-incidence matrix, so $Ap=f$ gives observed cell probabilities. Include normalization among the rows if needed. Let $B$ have first row $\mathbf1^{\mathsf T}$ and remaining rows $(H_j(i))_i$; let $b=(1,R_f\pi_1,\ldots)^{\mathsf T}$.

With no further cross-law structural restrictions,
\[
 \mathcal I(D)=\{p:Ap=f,\ p\geq\varepsilon\mathbf1\}
 \times\{q:Bq=b,\ q\geq\varepsilon\mathbf1\}. \tag{1}
\]
All laws in (1) have common full support, so equivalence imposes no additional linkage. An observed return distribution is not silently substituted for the full state distribution.

\section*{Identification by row spaces}
\begin{proposition}
Assume each factor of (1) has a feasible point with all coordinates strictly above $\varepsilon$. A linear physical functional $a^{\mathsf T}p$ is identified if and only if $a$ belongs to the row space of $A$. A pricing functional $h^{\mathsf T}q$ is identified if and only if $h$ belongs to the row space of $B$. Both laws are identified if and only if both matrices have column rank $n$.
\end{proposition}
\begin{proof}
If $a=A^{\mathsf T}v$, then $a^{\mathsf T}p=v^{\mathsf T}f$ is constant on the physical feasible set. If not, finite-dimensional orthogonality gives $z\in\ker A$ with $a^{\mathsf T}z\ne0$. At a strictly feasible $p_0$, both $p_0+tz$ and $p_0-tz$ remain feasible for sufficiently small $t>0$, and their targets differ. The pricing statement is identical. Applying the result to each coordinate vector yields the rank characterization.
\end{proof}

Strict feasibility matters: at boundary data, nonnegativity can uniquely determine a law even if the matrix is rank deficient. Also, if a structural model imposes $Q=P$ or a parametric pricing kernel, the Cartesian factorization (1) fails and the separate rank conditions need not be necessary.

For a payoff $H$ with vector $h$, the physical-versus-pricing expectation difference is $h^{\mathsf T}p-h^{\mathsf T}q$. Under (1) and strict feasibility it is identified exactly when $h$ lies in both row spaces: varying one factor while holding the other fixed proves necessity. For a traded payoff, the pricing expectation is observed, so only identification of its physical expectation remains. For an untraded payoff, both components can vary independently.

\section*{Sharp bounds, including likelihood ratios}
Linear target bounds are ordinary linear programs over (1), attained by compactness. For the difference above, write $\underline P_h,\overline P_h$ and $\underline Q_h,\overline Q_h$ for the separate LP bounds. Then the sharp interval is
\[
 [\underline P_h-\overline Q_h,\ \overline P_h-\underline Q_h].
\]
Independence of the two feasible factors makes both endpoints jointly attainable. This subtraction rule need not be sharp if a common structural parameter couples the factors.

For a state-specific likelihood ratio $m_i=q_i/p_i$, independent variation and positive denominators give exact bounds
\[
 \frac{\underline q_i}{\overline p_i}\leq m_i\leq
 \frac{\overline q_i}{\underline p_i}. \tag{2}
\]
Each endpoint is attainable by separately selecting its two extremizers. The vector identified set is not the Cartesian product of these coordinate intervals: probabilities sum to one and price constraints couple coordinates. In particular, $\sum_ip_im_i=1$ must hold. Reporting only (2) as the complete joint pricing-kernel set would be incorrect.

As a concrete example, take three states, observe only whether state 1 occurred, and observe only the bond price. Then $p_1=f$, $p_2+p_3=1-f$, while $q_1+q_2+q_3=1$. The physical mean of the indicator of state 1 is known, its pricing mean is not. Adding a traded claim paying exactly that indicator fixes $q_1$ but still leaves $p_2,p_3,q_2,q_3$ free. Adding a state-2 claim identifies all of $Q$ but does not identify the remaining physical split. Prices alone cannot repair missing historical observation without a cross-law model.

\section*{Equivalence and infinite spaces}
If the lower bound is removed but strict positivity retained, the feasible set is relatively open at support boundaries. Infima and suprema can then fail to be attained, and the upper likelihood-ratio bound can diverge as $p_i\downarrow0$. This is why equivalence is weaker than a uniform density-ratio bound.

For a standard Borel space, observing the full state law identifies $P$. If traded bounded payoffs form a measure-determining family, their correctly normalized expectations identify $Q$. When both measures are identified and $Q\ll P$, the Radon--Nikodym derivative is unique up to $P$-null sets. These are sufficient conditions, not a claim that every incomplete market identifies a kernel. One must distinguish uniqueness of a derivative for a given pair from identifying that pair through data.

\section*{An estimation guarantee and its conditioning}
Suppose independent historical observations yield empirical cell frequencies $\widehat f$, and prices have declared error intervals $|\widehat b_j-b_j|\leq e_j$. For $K$ observed cells, Hoeffding and a union bound imply
\[
 \Prb\left(\max_k|\widehat f_k-f_k|\leq
 \sqrt{\frac{\log(2K/\alpha)}{2N}}\right)\geq1-\alpha.
\]
Replacing equalities in (1) by these frequency bands and the price bands gives a set covering the true pair whenever the bands hold. Optimizing a target over that set yields simultaneous identification-and-sampling bounds. They need not be narrow.

If an identified square subsystem has smallest singular value $s>0$, inversion gives Euclidean estimation error bounded by data error divided by $s$. Weakly distinguishing payoff vectors can make $s$ arbitrarily small, so algebraic full rank alone does not imply a uniform estimation rate. A ratio estimate also requires a physical probability lower bound: when $p_i,\widehat p_i\geq\varepsilon$, its error is at most $|\widehat q_i-q_i|/\varepsilon+|q_i||\widehat p_i-p_i|/\varepsilon^2$.

\section*{Status and source}
The finite full-support benchmark is characterized exactly. General structural restrictions, infinite-dimensional completeness, dependence in historical samples and uniform estimation under weak identification remain unresolved.

Zhang Chen and Chen Kay (2026), \emph{Historical Developments in Probability Measures for Asset Pricing: From State Prices to Modern Pricing Kernels}, arXiv:2605.27658, Section 20. The primary PDF's discussion of physical/pricing-measure identification was inspected; the linear-algebra results above are proved here as restricted calculations. \url{https://arxiv.org/pdf/2605.27658}.

\end{document}
